\section{Deep Neural Networks \& Training}
\slideref{Session 3}

These notes accompany the third lecture. In Session~2 we established that a one-hidden-layer network with ReLU activations can approximate any continuous function --- but we stopped short of asking how to actually \emph{find} good parameters. This session addresses that question. We first generalise the architecture to arbitrary depth, then introduce the machinery needed to train it: a loss function to measure error, and gradient descent to minimise it.

% - - - - - - - - - - - - - - - - - --
\subsection{From Shallow to Deep Networks}
\slideref{Session 3: From Shallow to Deep, Matrix Notation}

The one-hidden-layer network from Session~2 computes, for input $\vx \in \R^d$:
\begin{align*}
h_j &= \act\!\left(\theta_{j0} + \sum_{i=1}^{d} \theta_{ji}\, x_i\right), \qquad j = 1, \ldots, k \\
\yhat &= \theta_0 + \sum_{j=1}^{k} \theta_j\, h_j.
\end{align*}
The natural extension is to stack additional layers: the output of one layer becomes the input of the next. Each layer applies a linear transformation followed by a nonlinear activation. Stacking layers allows the network to learn increasingly abstract representations --- the first layer might detect simple patterns in the input, the second layer combinations of those patterns, and so on.

To handle multiple layers cleanly, scalar notation becomes cumbersome. We instead collect the weights of layer $l$ into a matrix $\thetamat{l}$ and the biases into a vector $\thetavec{l}_0$. For a single hidden layer with $k$ neurons and $d$-dimensional input:
\[
\boldsymbol{\Theta} =
\begin{pmatrix}
\theta_{11} & \cdots & \theta_{1d} \\
\vdots & \ddots & \vdots \\
\theta_{k1} & \cdots & \theta_{kd}
\end{pmatrix}
\in \R^{k \times d}, \qquad
\thetavec{}_0 =
\begin{pmatrix} \theta_{10} \\ \vdots \\ \theta_{k0} \end{pmatrix}
\in \R^k.
\]
The hidden layer computation then becomes:
\[
\vh = \act\!\left(\thetavec{}_0 + \boldsymbol{\Theta}\vx\right) \in \R^k,
\]
where $\act$ is applied elementwise. Each row of $\boldsymbol{\Theta}$ contains the weights of one neuron; each column contains the weights from one input dimension.

% - - - - - - - - - - - - - - - - - --
\subsection{Layer Notation and Function Composition}
\slideref{Session 3: Layer Notation, Deep Network as Function Composition, Dimensions}

We index layers $l = 0, 1, \ldots, L$, where $l=0$ denotes the input and $l=L$ the output. The forward pass is defined recursively:
\[
\vh^{(0)} = \vx, \qquad \vh^{(l)} = \act\!\left(\thetavec{l}_0 + \thetamat{l}\, \vh^{(l-1)}\right), \quad l = 1, \ldots, L-1,
\]
with output
\[
\yhat = \thetavec{L}_0 + \thetamat{L}\, \vh^{(L-1)}.
\]
Note that the output layer applies no activation function --- the raw linear output is passed directly as the prediction. This is the standard convention for regression; we revisit output activations when we discuss classification.

The weight matrix at layer $l$ has shape $\thetamat{l} \in \R^{k_l \times k_{l-1}}$, where $k_l$ is the width of layer $l$ and $k_0 = d$ is the input dimension. The rows correspond to neurons in layer $l$; the columns correspond to neurons (or inputs) in layer $l-1$.

Viewed as a whole, the deep network is a composition of functions:
\[
\yhat = f_\theta(\vx) = \left(h^{(L)}_\theta \circ \cdots \circ h^{(1)}_\theta\right)(\vx) = \thetavec{L}_0 + \thetamat{L}\,\act\!\left(\cdots\,\act\!\left(\thetavec{1}_0 + \thetamat{1}\vx\right)\cdots\right).
\]
The total number of parameters is the sum over all layers of the number of weights plus biases:
\[
|\pvec| = \sum_{l=1}^{L} k_l \cdot (k_{l-1} + 1).
\]
For example, a network with $d=3$ inputs, two hidden layers of width $k_1 = k_2 = 4$, and $m=2$ outputs has:
\[
4 \cdot (3+1) + 4 \cdot (4+1) + 2 \cdot (4+1) = 16 + 20 + 10 = 46 \text{ parameters.}
\]
Real networks used in practice have millions or billions of parameters. Finding good values for all of them jointly from data is the central challenge of deep learning.

% - - - - - - - - - - - - - - - - - --
\subsection{Batch Formulation}
\slideref{Session 3: Batch Formulation}

In practice, we process multiple training samples simultaneously in a \textbf{batch}. Given $N$ samples, we stack the inputs as rows of a matrix $\mathbf{X} \in \R^{N \times d}$, and the hidden activations of layer $l$ form a matrix $\mathbf{H}^{(l)} \in \R^{N \times k_l}$.

The batch forward pass for layer $l$ is:
\[
\mathbf{H}^{(l)} = \act\!\left(\mathbf{1}\thetavec{l}_0{}^\top + \mathbf{H}^{(l-1)} \left(\thetamat{l}\right)^\top\right),
\]
where $\mathbf{1} \in \R^N$ is a vector of ones that broadcasts the bias across all samples. This is precisely what PyTorch computes when we write \texttt{H = X @ W.T + b}: the batch dimension comes first, and the matrix multiply operates across the feature dimension.

Processing data in batches has two benefits. First, it is computationally efficient: modern hardware (GPUs, TPUs) is designed for large parallel matrix operations, and a single batched matrix multiply is far faster than $N$ sequential vector operations. Second, it provides more stable gradient estimates than processing one sample at a time.

% - - - - - - - - - - - - - - - - - --
\subsection{The Learning Problem}
\slideref{Session 3: The Learning Problem}

We now have a network $f_\theta$ and a dataset $\dataset = \{(\vx^{(i)}, y^{(i)})\}_{i=1}^{\nsamples}$. The goal is to find parameters $\pvec$ such that $f_\theta(\vx) \approx y$ on unseen data. Two questions immediately arise:

\begin{enumerate}
    \item How do we \textbf{measure} how wrong $f_\theta$ is? $\Rightarrow$ \emph{loss function}
    \item How do we \textbf{find} good $\pvec$ systematically? $\Rightarrow$ \emph{optimisation}
\end{enumerate}

As illustrated in Figure~\ref{fig:learning_problem}, many different functions could plausibly fit the data. We need a principled criterion to distinguish better from worse.

\begin{figure}[h]
    \centering
    \resizebox{0.5\textwidth}{!}{% Common/tikz/learning_problem.tex
\begin{tikzpicture}[scale=0.8]
\tikzstyle{pt}    = [circle, draw, fill=thwsblue!40, minimum size=0.1cm, inner sep=0pt]
\tikzstyle{annot} = [font=\scriptsize]

% Axes
\draw[-Stealth, thick] (-0.3, 0) -- (3.8, 0) node[right, annot] {$x$};
\draw[-Stealth, thick] (0, -0.3) -- (0, 3.5) node[above, annot] {$y$};

% Data points
\node[pt] at (0.3, 0.55) {};
\node[pt] at (0.6, 0.7) {};
\node[pt] at (1.0, 0.60) {};
\node[pt] at (1.5, 0.8) {};
\node[pt] at (1.8, 1.10) {};
\node[pt] at (2.0, 1.0) {};
\node[pt] at (2.2, 1.45) {};
\node[pt] at (2.45, 1.90) {};
\node[pt] at (2.5, 2.30) {};
\node[pt] at (2.8, 2.60) {};
\node[pt] at (3.0, 3.30) {};

% Fit 1: one breakpoint
\draw[thick, thwspetrol]
    (-0.05, 0.1) -- (2.0, 0.7) -- (3.5, 2.8)
    node[annot, right, anchor=west] {$f^a_\theta$};

% Fit 2: two breakpoints
\draw[thick, thwsorange]
    (-0.05, 0.4) -- (1.2, 0.65) -- (2.3, 1.55) -- (3.5, 3.2)
    node[annot, right, anchor=west] {$f^b_\theta$};

% Fit 3: three breakpoints
\draw[thick, carnelian]
    (-0.05, 0.55) -- (1.0, 0.70) -- (1.9, 1.15) -- (2.5, 2.40) -- (3.5, 3.5)
    node[annot, right, anchor=west] {$f^c_\theta$};

% % Question
% \node[annot, font=\normalsize] at (1.8, -0.7) {Which $f_\theta$ is best?};

\end{tikzpicture}}
    \caption{Multiple piecewise linear functions fit the training data to varying degrees. A loss function provides a scalar measure of fit quality.}
    \label{fig:learning_problem}
\end{figure}

% - - - - - - - - - - - - - - - - - --
\subsection{Loss Functions}
\slideref{Session 3: What is a Loss Function?, Regression: Mean Squared Error}

A \textbf{loss function} $\ell(\yhat, y)$ measures the discrepancy between a prediction $\yhat$ and a true label $y$ for a single sample. It satisfies $\ell(\yhat, y) = 0$ when $\yhat = y$ and $\ell(\yhat, y) > 0$ otherwise, increasing as the prediction worsens. The \textbf{dataset loss} averages the per-sample loss over all training examples:
\[
\loss(\pvec) = \frac{1}{\nsamples} \sum_{i=1}^{\nsamples} \ell\!\left(f_\theta(\vx^{(i)}), y^{(i)}\right).
\]
The choice of loss function depends on the task. For regression we use Mean Squared Error; for binary classification, Binary Cross-Entropy; for multiclass classification, Cross-Entropy with a softmax output.

\subsubsection{Mean Squared Error}

For regression with scalar output $y \in \R$ and prediction $\yhat \in \R$, the per-sample loss is:
\[
\ell(\yhat, y) = (\yhat - y)^2.
\]
The dataset loss is accordingly:
\[
\loss(\pvec) = \frac{1}{\nsamples}\sum_{i=1}^{\nsamples}\!\left(\yhat^{(i)} - y^{(i)}\right)^2.
\]
Three properties make squaring the natural choice. First, prediction errors can be positive or negative and would cancel if averaged directly --- squaring removes the sign. Second, squaring penalises large errors more heavily than small ones, which is usually desirable. Third, the squared function is differentiable everywhere, which is essential for gradient-based optimisation.

The output layer for regression has no activation function: $\yhat = \thetavec{L}_0 + \thetamat{L}\vh^{(L-1)} \in \R$. The raw linear output is used directly as the prediction.

\begin{figure}[h]
    \centering
    \resizebox{0.5\textwidth}{!}{% Common/tikz/mse_loss.tex
\begin{tikzpicture}[scale=0.8]
\tikzstyle{pt}    = [circle, draw, fill=thwsblue!40, minimum size=0.1cm, inner sep=0pt]
\tikzstyle{annot} = [font=\scriptsize]

% Axes
\draw[-Stealth, thick] (-0.3, 0) -- (3.8, 0) node[right, annot] {$x$};
\draw[-Stealth, thick] (0, -0.3) -- (0, 3.5) node[above, annot] {$y$};

% Piecewise linear fit
\draw[thick, carnelian]
    (-0.05, 0.4) -- (1.0, 0.65) -- (1.9, 1.15) -- (2.5, 2.30) -- (3.5, 3.4)
    node[annot, right, anchor=west] {$f_\theta$};

% Data points and residuals
\foreach \x/\y/\yfit in {
    0.3/0.55/0.48,
    0.6/0.7/0.54,
    1.0/0.60/0.65,
    1.5/0.8/0.90,
    1.8/1.10/1.10,
    2.0/1.0/1.35,
    2.2/1.45/1.7,
    2.5/2.30/2.30,
    2.8/2.60/2.6,
    3.0/3.30/2.85} {
    % Residual line
    \draw[thwspetrol!60, thin] (\x, \y) -- (\x, \yfit);
    % Data point
    \node[pt] at (\x, \y) {};
}

% Residual label
\draw[thwspetrol!60, -Stealth, thin] (1.35, 1.3) -- (1.9, 1.18);
\node[annot, thwspetrol] at (1.1, 1.5) {$\yhat^{(i)}\!-y^{(i)}$};

\end{tikzpicture}}
    \caption{Mean squared error illustrated geometrically. Vertical lines show the residuals $\yhat^{(i)} - y^{(i)}$ between the predicted function and the data points. MSE averages the squares of these residuals.}
    \label{fig:mse_loss}
\end{figure}

% - - - - - - - - - - - - - - - - - --
\subsection{Gradient Descent}
\slideref{Session 3: The Optimization Problem, Gradient Descent, The Key Question}

Training the network means finding parameters that minimise the dataset loss:
\[
\pvec^* = \arg\min_{\pvec}\ \loss(\pvec).
\]
Unlike linear regression, there is no closed-form solution: $\loss(\pvec)$ is nonlinear in $\pvec$, the parameter space has millions of dimensions, and the loss surface is non-convex with many local minima. We therefore resort to iterative, gradient-based optimisation.

The \textbf{gradient} $\nabla_{\pvec} \loss(\pvec)$ is the vector of partial derivatives of $\loss$ with respect to each parameter:
\[
\nabla_{\pvec} \loss(\pvec) = \frac{\partial \loss}{\partial \pvec} = \left(\frac{\partial \loss}{\partial \theta_1}, \frac{\partial \loss}{\partial \theta_2}, \ldots\right).
\]
It points in the direction of steepest ascent of the loss surface. To minimise the loss we step in the opposite direction. This gives the \textbf{gradient descent} update rule:
\[
\pvec \leftarrow \pvec - \eta\, \nabla_{\pvec} \loss(\pvec),
\]
where $\eta > 0$ is the \textbf{learning rate}, controlling the step size. We repeat this update until convergence.

\begin{figure}[h]
    \centering
    \resizebox{0.5\textwidth}{!}{% PLACEHOLDER: gradient_descent_2d
\begin{tikzpicture}
  \draw[thick, thwsgrey, rounded corners=4pt] (0,0) rectangle (4,2.5);
  \node[thwsgrey, font=\scriptsize\ttfamily] at (2,1.25) {gradient\_descent\_2d};
\end{tikzpicture}
}
    \caption{Gradient descent on a quadratic loss. Starting from $\theta_0$, successive steps approach the minimum $\theta^*$. The step size is proportional to the learning rate $\eta$.}
    \label{fig:gradient_descent}
\end{figure}

The learning rate is a critical hyperparameter. Too large and the updates overshoot, potentially diverging. Too small and convergence is impractically slow. Choosing a good learning rate --- and adapting it during training --- is a topic we return to in the dedicated optimisation session.

The central computational challenge is evaluating $\nabla_{\pvec} \loss(\pvec)$ efficiently. The loss $\loss(\pvec)$ is a composition of many nested functions --- one per layer of the network. Differentiating through this composition na\"{i}vely would be prohibitively expensive. The answer is \textbf{backpropagation}, which we develop in the next session.
